\section{Discussion}
\paragraph{Where the model earns its keep.} When a relative in the same genus has a measured growth rate, simply taking the
median of the relatives gets close to the model, which then adds mostly a better ranking. The model pulls ahead as the
relatives get further away, and that is where most environmental genomes sit: \GEMShareNoGenus\% of the high-quality GEM
genomes have no labeled relative in their genus and \GEMShareNewOrder\% none in their order. In that regime the accuracy comes
from the genes, through gRodon's codon statistics and the gene features, while the tree position and the 16S gene add
little. We would use the predictions to compare and rank genomes rather than to quote a single doubling time, and for
comparisons across habitats we would use the doubling time at \TRefC\,\textdegree C, which removes the part of the difference
that temperature explains. Slow growth remains the weak point. Slow growers are ranked well but rarely predicted above a day,
so the \GEMNSlow slow GEM genomes are a short list rather than a count, and the slow-growth probability is the better guide.
The clearest next gain is to add protein language-model embeddings to the features.

\paragraph{Labels and inputs.} The \NExcludedKevbrinaWord values from a single study were enough to produce a \KevFoldWith-fold effect, significant
with relatedness controlled, that would have read as support for an attractive biological idea. The values agreed with each
other, so no comparison with relatives could have flagged them; only tracing them back to their source did. We would check
values from studies that contribute several species, or that differ from their source record, before using them. The 16S
result is the same trap on the input side: copies retrieved by name and copies cut from the genome were similar enough to pass
our drift check and different enough to cost accuracy. Finally, the measured temperature lowers gRodon's error from
\RelGRodonNoTempErr- to \RelGRodonErr-fold, an advantage an uncultured genome does not have. Predictors meant
for metagenome-assembled genomes should be judged without it, and with whole lineages held out.

\paragraph{Slow growth and energy.} Organisms in energy-poor environments are often thought to grow slowly because their
energy supply is thin \citep{hoehler2013,lever2015}. Across species, several energy-related genes do track doubling time, but
none survives control for relatedness and for the number of inputs tested, in the main analysis or in the three checks of
Supplementary Table~\ref{S-tab:pgls}. These associations follow the tree. What remains is that slow growers carry fewer
transporter genes, in line with what is known about oligotrophic and copiotrophic marine bacteria \citep{lauro2009}. Our
labeled organisms all grow in culture and few live on the thinnest energy budgets, so this argues against the simplest form
of the low-energy explanation among cultured organisms, not against it in general.

\paragraph{Limitations.} The labels are maximal rates in culture, which can be far above rates in the environment
\citep{gibson2018,brown2016}. The training set is small (\NGenomes genomes, \NSlow slower than a day), and predictions shrink
toward the middle of the range. Some labels are surely still wrong after the audit. The GEM taxonomy comes from an older GTDB
release, so the novelty levels we report are worst cases. On species outside both compilations the ranges held fewer values
than they should, and missed most slow growers, so for a new lineage the true uncertainty is larger than the range suggests. In our tests, adding labeled species helped more than changing the learner, so new growth measurements for lineages
without cultured members would probably improve the predictions most.
